\documentclass[11pt]{article}
\usepackage[margin=1in]{geometry}
\usepackage{booktabs}
\usepackage{tabularx}
\usepackage{enumitem}
\usepackage[hidelinks]{hyperref}
\usepackage{array}
\usepackage{float}
\newcolumntype{L}{>{\raggedright\arraybackslash}X}
\newcolumntype{R}{>{\raggedleft\arraybackslash}X}
\setlist{nosep, leftmargin=*}

\title{\textbf{CS622 Project: Progress Since the Mid-Review}\\[4pt]
\large L2 TLB Support for Sparse Huge Pages}
\date{}

\begin{document}
\maketitle

\noindent Since the mid-review we fixed correctness bugs in our L2 TLB designs, made CoLT-FA work under a real operating system, and evaluated all six designs in gem5 full-system mode. Every SPEC CPU2017 benchmark that can run in this setup passes on every design, and the remaining failures are independent of the TLB. On graph workloads, CoLT-FA cuts page walks by 24--66\% relative to the NINE L2 it is built on.

\section{At a Glance}
\begin{table}[H]
\centering
\begin{tabularx}{\textwidth}{@{}l L@{}}
\toprule
\textbf{Workstream} & \textbf{Result} \\
\midrule
Correctness fixes & Six designs fixed: Inclusive, Exclusive, NINE, Agile/NaiveFP, Perforated Pages, CoLT-FA \\
CoLT-FA in full-system mode & Coalescing moved into the hardware page walker; 0 oracle mismatches \\
SE-mode verification & 0 oracle mismatches on all designs; remaining failures identical on stock gem5 \\
Full-system SPEC CPU2017 & 37 of 43 benchmarks pass on every design; all failures are TLB-independent \\
TLB-intensive workloads & GAPBS (6 kernels), memcached, XSBench, STREAM, SPEC reference inputs \\
CoLT-FA vs.\ NINE & $-$24 to $-$66\% page walks on GAPBS; $-$44\% on xalancbmk (ref) \\
\bottomrule
\end{tabularx}
\end{table}

\section{Correctness Fixes}
An audit after the mid-review showed that the designs were correct for single-process SE runs but not under a real operating system.

\begin{table}[H]
\centering
\begin{tabularx}{\textwidth}{@{}l L@{}}
\toprule
\textbf{Design} & \textbf{Fixes} \\
\midrule
All designs & L2 invalidated on page unmap and CR3 flush; ASID-aware L2 removal; L2 state saved in gem5 checkpoints; L2 statistics; geometry, hash and replacement made configurable \\
Exclusive & Promotion left duplicate entries under SMT: 391{,}950 inclusion violations reduced to 0 \\
Perforated Pages & Six bugs fixed, including hard-coded permissions on L2 hits \\
Agile / NaiveFP & Prefetcher generalised into a parameterised design \\
CoLT-FA & Rebuilt on the fixed NINE L2; ASID-tagged entries; real permissions from the page walk; invalidated on flushes, unmaps and checkpoint restore; coalescing moved into the page walker (Section~\ref{sec:colt}) \\
\bottomrule
\end{tabularx}
\end{table}

\section{CoLT-FA in Full-System Mode}
\label{sec:colt}
The mid-review CoLT-FA buffer formed its coalesced entries from gem5's emulated page table, which exists only in SE mode, so in full-system mode it never coalesced anything. We rebuilt it on the fixed NINE L2 TLB, the baseline used for CoLT in the mid-review. As in the CoLT paper, coalescing now happens in the hardware page walker: when a walk fetches a 4\,KB leaf PTE, the same 64-byte cache line holds seven neighbouring PTEs, and the TLB inserts the longest run of present, physically contiguous PTEs with identical permissions (up to 8 pages) into an 8-entry fully associative buffer. The buffer is probed after an L1 miss and before the L2. A run-time switch disables it, so the same binary also serves as the NINE baseline for matched comparisons.

\section{Verification in SE Mode}
An oracle checks every translation against the page table, together with an inclusion-invariant checker. No design shows an oracle mismatch.

\begin{table}[H]
\centering
\begin{tabularx}{\textwidth}{@{}l L@{}}
\toprule
\textbf{Design} & \textbf{Result} \\
\midrule
Inclusive & 0 mismatches; SPEC test inputs: 31/31 runnable invocations pass, output byte-identical to stock gem5 \\
Exclusive & 94/120 stress cases pass; all 26 failures also fail on stock gem5 \\
NINE & 94/120 pass; 0 design-specific failures; 12 configurations swept clean \\
Perforated Pages & 36/44 pass, the other 8 identical on stock gem5; 31 configurations, 0 mismatches \\
Agile / NaiveFP & 16+ prefetch configurations measured \\
CoLT-FA & 36/36 runs pass (Atomic and Timing CPUs, CoLT on and off); 0 mismatches, 0 stale entries \\
\bottomrule
\end{tabularx}
\end{table}

\noindent The stock-gem5 failures come from SE-mode limits (\texttt{setitimer}, \texttt{MAP\_SHARED}, multi-process remaps).

\section{Full-System SPEC CPU2017}
All 43 SPEC CPU2017 benchmarks (test inputs) were run in gem5 full-system mode on an Ubuntu 24.04 guest (Linux 6.8) for each of the six designs, about 300 simulations across the department lab machines. Each benchmark runs on gem5's atomic CPU with the TLB under test, and its output is validated against SPEC's reference output. Each also gets a 200M-instruction window on the detailed timing CPU for performance statistics.

\begin{table}[H]
\centering
\begin{tabularx}{\textwidth}{@{}l c L@{}}
\toprule
\textbf{Design} & \textbf{PASS / 43} & \textbf{Failures} \\
\midrule
Inclusive, Exclusive, NINE, Perforated & 37 & perlbench\_r, perlbench\_s, deepsjeng\_s \\
Agile, CoLT-FA & 36 (+1 running) & perlbench\_r, perlbench\_s, deepsjeng\_s \\
\bottomrule
\end{tabularx}
\caption{Full-system SPEC CPU2017 results (test inputs).}
\end{table}

\noindent Every failure is independent of the TLB design:
\begin{itemize}
  \item \textbf{perlbench\_r, perlbench\_s:} fail on every design due to a gem5 x87 floating-point emulation bug (long-double formatting). Both pass natively with the identical SPEC harness.
  \item \textbf{631.deepsjeng\_s:} needs 6.9\,GB, but gem5's x86 full-system board caps the guest at 3\,GB. It passes in SE mode with 8\,GB on all six designs.
  \item \textbf{619.lbm\_s} (3.2\,GB, same limit) has no full-system result; it passes in SE mode with 8\,GB on all six designs. \textbf{525.x264\_r} and \textbf{625.x264\_s} ($\sim$1 trillion instructions each) cannot finish in simulation and are reported with timing-window statistics only.
\end{itemize}

\section{Performance on TLB-Intensive Workloads}
SPEC's test inputs barely stress the TLB (at most 1.1 page walks per 1000 instructions), so we added the workloads suggested for this study: the GAPBS graph kernels, the memcached key-value store, XSBench, STREAM, and SPEC reference inputs. Each is measured on the timing CPU after an atomic warm-up past its set-up phase. The metric is page walks per 1000 instructions (lower is better).

\subsection{CoLT-FA vs.\ NINE (matched)}
Both columns use the same binary with CoLT switched off or on, so the comparison isolates CoLT.

\begin{table}[H]
\centering
\begin{tabularx}{\textwidth}{@{}L R R R@{}}
\toprule
\textbf{Workload (4\,KB pages only)} & \textbf{NINE} & \textbf{CoLT-FA} & \textbf{Change} \\
\midrule
GAPBS bfs (Kronecker $2^{20}$) & 1.23 & 0.42 & $-$66\% \\
GAPBS cc (uniform $2^{20}$) & 0.46 & 0.16 & $-$66\% \\
GAPBS sssp & 6.35 & 2.58 & $-$59\% \\
GAPBS bc & 26.55 & 16.91 & $-$36\% \\
GAPBS pr & 54.76 & 38.42 & $-$30\% \\
GAPBS tc & 3.25 & 3.22 & $-$1\% \\
XSBench & 14.26 & 12.29 & $-$14\% \\
memcached (random GETs) & 1.18 & 1.15 & $-$3\% \\
STREAM & 0.40 & 0.40 & 0\% \\
\bottomrule
\end{tabularx}
\caption{Page walks per 1000 instructions with transparent huge pages disabled.}
\end{table}

\noindent With transparent huge pages enabled, CoLT-FA still reduces page walks on GAPBS bfs ($-$50\%), pr ($-$51\%), bc ($-$29\%) and sssp ($-$24\%). On SPEC reference inputs it reduces them on xalancbmk ($-$44\%), omnetpp ($-$22\%) and lbm ($-$78\%), and leaves mcf unchanged. CoLT gains most when the operating system leaves 4\,KB pages physically contiguous, and least on random-access workloads (memcached, XSBench) and streaming (STREAM).

\subsection{Across Designs}
\begin{table}[H]
\centering
\begin{tabularx}{\textwidth}{@{}L R R R R@{}}
\toprule
\textbf{Workload (4\,KB pages only)} & \textbf{Inclusive} & \textbf{Exclusive} & \textbf{Agile} & \textbf{CoLT-FA} \\
\midrule
GAPBS bfs & 1.25 & 1.22 & 1.24 & 0.42 \\
GAPBS pr & 55.05 & 51.82 & 54.69 & 38.42 \\
GAPBS bc & 26.74 & 25.85 & 26.65 & 16.91 \\
XSBench & 14.39 & 13.97 & 14.25 & 12.29 \\
\bottomrule
\end{tabularx}
\caption{Page walks per 1000 instructions, transparent huge pages disabled.}
\end{table}

\noindent Exclusive is the best inclusion policy (up to 6\% fewer walks than Inclusive, because it does not duplicate entries), and the Agile prefetcher gives little benefit on these workloads. On SPEC reference inputs with huge pages enabled, mcf (3.4--3.7) and omnetpp (2.2--2.3) are the most TLB-intensive benchmarks, and the inclusion policies differ by under 10\%.

\subsection{A Caveat on Perforated Pages}
As implemented, the Perforated Pages L2 stores a separate physical address for each of the 512 4\,KB pages in a 2\,MB region, with no check that they are physically contiguous. This gives it roughly 512 times the reach of the other designs, which real hardware cannot provide, and it shows near-zero page walks on most workloads (e.g.\ mcf ref 0.12 vs.\ 3.62 for NINE). With the paper's rule enforced (one physical base per region, holes allowed), it performs worse than NINE on most workloads (mcf ref 10.35; GAPBS pr 88.6 vs.\ 54.8) and better only on XSBench (4.94 vs.\ 14.26). We therefore treat CoLT-FA as our main realistic design.

\section{Limitations}
\begin{itemize}
  \item gem5's x86 full-system board limits the guest to 3\,GB, which excludes deepsjeng\_s and lbm\_s from full-system runs (both pass in SE mode).
  \item perlbench is affected by a gem5 x87 emulation bug, independent of the TLB.
  \item Performance is measured over sampled timing windows (200M--1.5B instructions), not whole runs.
  \item With huge pages enabled, the guest's physical memory layout depends on the disk image, so designs are compared only within runs on the same machine.
  \item Redis runs did not complete under gem5; memcached covers the key-value store case.
\end{itemize}

\section{Conclusion}
After making the six L2 TLB designs correct under a real operating system, all of them pass every SPEC CPU2017 benchmark that the simulator can run. Among the realistic designs, CoLT-FA gives the largest reduction in page walks, up to 66\% on graph workloads, and it remains useful even when transparent huge pages are available.

\end{document}
